\chapter{The Density Cardinal}

\section{Introduction}

For $A\subseteq\omega$, let $d_n(A)=\frac{\lvert A\cap n\rvert}{n}$ and
\begin{align*}
    \underline{d}(A) & = \liminf_{n\to\infty} d_n(A)\\
    \ol{d}(A) & = \limsup_{n\to\infty} d_n(A)
\end{align*}
If $\underline{d}(A)=\ol{d}(A)$ we call the common value $d(A)$ the \lasereye{(asymptotic) density} of $A$. Note that, when $A$ has asymptotic density, we have that $d(A)\in[0,1]$. Note that, analogously to Riemann's Rearrangement Theorem, we can prove the following theorem:

\begin{lemma}{Riemann for Density}{riemann_density}
    Given an infinite and coinfinite set $A$, we have
    \begin{itemize}
        \item For every $r\in[0,1]$, there is a $p\in\sym$ such that $d(p[A])=r$;
        \item There is a $p\in\sym$ such that $p[A]$ does not have asymptotic density
    \end{itemize}
\end{lemma}

\begin{proof}
    Given $r\in(0,1)$, let $p_1=1$ and $B_1=\{1\}$. Now, let $q_1$ be the least $N>p_1$ such that $d_N(B_1)<r$. Now, define $p_i,B_i$ and $q_i$, for $i>1$ recursively as follows:

    Let $p_n$ be the least $N>q_{n-1}$ such that $B_n=B_{n-1}\cup\{q_{n-1}+1,q_{n-1}+2,\dots,N\}$ has $d_N(B_n)>r$. And $q_n$ be the least $N>p_n$ such that $d_N(B_n)<r$.

    Now, let $B=\bigcup_{i\geq1}B_i=\{b_1,b_2,b_3,\dots\}$ with $b_1<b_2<\dots$ and $B^c=\{\tilde{b_1},\tilde{b_2},\dots\}$ also in order. Analogously, let $A=\{a_1,a_2,\dots\}$ and $A^c=\{\tilde{a_1},\tilde{a_2},\dots\}$ both in order. We can then define $p\in\sym$ such that $p(a_i)=b_i$ and $p(\tilde{a_i})=\tilde{b_i}$, which is obviously a permutation in $\omega$.

    With this we have that, since $p_i$ is the least $N$ with $d_N(p[A])>r$, for any $i>1$
    $$d_{p_i-1}(p[A])<r<d_{p_i}(p[A])$$
    therefore
    \begin{align*}
        0<d_{p_i}(p[A])-r & < d_{p_i}(p[A])-d_{p_i-1}(p[A])\\
        & = \frac{\lvert p[A]\cap p_i\rvert}{p_i}-\frac{\lvert p[A]\cap (p_i-1)\rvert}{p_i-1}\\
        & < \frac{\lvert p[A]\cap p_i\rvert-\lvert p[A]\cap (p_i-1)\rvert}{p_i-1}
    \end{align*}
    and, since the numerator of the RHS is either $0$ or $1$, and $p_i-1\to\infty$, then, when $i\to\infty$, the RHS converges to $0$, therefore $d_{p_i}(p[A])\to r$. Similarly, since $d_{q_i}(p[A])<r<d_{q_i-1}(p[A])$, we have that $d_{q_i}(p[A])\to r$. And, because
    $$d_{q_i}(p[A])\leq d_{q_i-1}(p[A])\leq \dots\leq d_{p_i-2}(p[A])\leq d_{p_i-1}(p[A])\leq d_{p_i}(p[A])$$
    then $d_n(p[A])\to r$.

    The cases for $p\in\sym$ s.t. $d(p[A])$ is $0$ or $1$, and where $p[A]$ does not have asymptotic density are analogous.
\end{proof}

\begin{lemma}{(\cite{brech})}{aux_dd_uncountable}
    Let $(a_n)$ be a sequence in $\omega$ satisfying that there exists $n_0\in\omega$ such that for all $n\geq n_0$, we have that $a_{n+1}-a_n>2^n$. Then $A=\{a_n\}_{n\geq0}$ is infinite-coinfinite and $d(A)=0$.
\end{lemma}

\begin{proof}
    It's clear that $A$ is infinite-coinfinite. Now, since $(a_n)$ is increasing for $n\geq n_0$, we have that
    $$a_{n+1}>a_n+2^n>a_{n_0}+2^n$$
    therefore, if $k\in\omega$ is s.t. $a_{n_0}+2^{n-1}\leq k\leq a_{n_0}+2^n<a_{n+1}$, then $A\cap k\subseteq\{a_i:i\leq n\}$, for $n\geq n_0$, and we will have that
    $$d_k(A)=\frac{\lvert A\cap k\rvert}{k}\leq\frac{n+1}{k}\leq\frac{n+1}{a_{n_0}+2^{n-1}}$$
    which converges to $0$ when $k\to\infty$, therefore $d(A)=0$.
\end{proof}

\begin{theorem}{(\cite{brech})}{dd_uncountable}
    \begin{center}
        $\mf{dd}$ is uncountable.
    \end{center}
\end{theorem}

\begin{proof}
    Let $P=\{p_n:n\in\omega\}\subseteq\sym$ and assume WLOG that $p_0=\mathrm{id}$. Let $a_0=1$ and, recursively, if $a_n$ is already defined, let $a_{n+1}>a_n$ be such that $p_i(a_{n+1})\geq p_i(a_n)+2^n$, for all $i\leq n$, which can be done, since $p_i(k)\to\infty$. Then let $A=\{a_n:n\in\omega\}$. By the previous lemma $d(p[A])=d(A)=0$, for all $p\in P$.
\end{proof}

\section{Characterization of \texorpdfstring{$\mf{dd}$}{dd}}

The main theorem of this subsection is the characterization $\mf{dd}=\textbf{non}(\mc{M})$. To prove this using GT, we will find relational systems $\mc{A}_1$, $\mc{A}_2$ and $\mc{B}$ such that we have
$$\mc{A}_1\leq_T\mc{B}\leq_T\mc{A}_2\qquad\mf{d}(\mc{A}_1)=\mf{d}(\mc{A}_2)=\textbf{non}(\mc{M})\qquad\mf{d}(\mc{B})=\mf{dd}$$

It's the usual that $\mc{A}_2=(\mc{M},2^\omega,\not\ni)$ works, where $\mc{M}$ are the meager sets of $2^\omega$. And, for $\mf{dd}$, we can take $\mc{B}=(D,\sym,R)$, where $D$ are the infinite-coinfinite substets of $\omega$ with density, and $xRp$ if $d(x)\neq d(p(x))$, as expected.

We will first prove\footnote{Here, again, we will use that $\sym$ is homeomorphic to $\omega^\omega$ in $2^\omega$} the easier inequality $\mc{B}\leq_T\mc{A}_2$, you should try to find a reduction for $\mc{A}_2\leq_T\mc{B}$, you will see that's not too easy (in fact it's impossible), this is why we will later find other relational system $\mc{A}_1$ to work with the same dominating number.

\begin{lemma}{(\cite{brech})}{dd_non}
    $$(D,\sym,R)\leq_T(\mc{M},2^\omega,\not\ni)$$
\end{lemma}

\begin{proof}
    We need to define $\varphi_-:D\to\mc{M}$ and $\varphi_+:\sym\to\sym$ such that, for all $A\in D$ and $p\in\sym$, if $p\notin\varphi_-(A)$, then $d(A)\neq d(\varphi_+(p)(A))$. Lets take $\varphi_+=\text{id}$ as simple as possible, so we need $\varphi_-(A)$ such that, if $p\notin\varphi_-(A)$, then $d(A)\neq d(p(A))$.

    We can't choose $p$ such that $d(p(A))=1$ (as in $\sum_na_{p(n)}=\infty$), since $\mf{dd}$ is the analogous cardinal of $\mf{rr}'$ and not $\mf{rr}$. Then we will choose another approach:

    Lets try to force $d(p(A))$ diverge by oscillation, so $d_{n_k}(p(A))\to 0$ and $d_{m_k}(p(A))\to1$ for two different subsequences. Then, $\forall\varepsilon>0$ and $\ell\in\omega$, there are $n_0,n_1\geq\ell$ such that $d_{n_0}(p(A))<\varepsilon$ and $d_{n_0}(p(A))>1-\varepsilon$. Then, a natural choice for $\varphi_-(A)^c$ would be
    \begin{align*}
        \varphi_-(A)^c & := \left\{p\in\sym:\forall\varepsilon>0,\forall\ell\in\omega,\exists n_0,n_1\geq\ell,d_{n_0}(p(A))<\varepsilon\text{ and }d_{n_1}(p(A))>1-\varepsilon\right\}\\
        & = \left\{p\in\sym:\forall k,\ell\in\omega,\exists n_0,n_1\geq\ell,d_{n_0}(p(A))<\frac1k\text{ and }d_{n_1}(p(A))>1-\frac1k\right\}\\
        & = \bigcap_{k,\ell\in\omega}\left\{p\in\sym:\exists n_0,n_1\geq\ell,d_{n_0}(p(A))<\frac1k\text{ and }d_{n_1}(p(A))>1-\frac1k\right\}\\
        & = \bigcap_{k,\ell\in\omega}A_{k,\ell}
    \end{align*}
    as we saw in the proof of \cref{th:rr_non}, to show that $\varphi_-(A)\in\mc{M}$, it's sufficient to prove that $\varphi_-(A)^c$ is the intersection of countably many dense sets, so we will show that each $A_{k,\ell}$ is open and dense. If $p\in A_{k,\ell}$, lets $m=\max\{n_0,n_1\}$, then, for $q\in\mc{N}(p\restriction_{m+1})$, clearly $d_{n_0}(q(A))=d_{n_0}(p(A))<\frac1k$ and $d_{n_1}(q(A))=d_{n_1}(p(A))>1-\frac1k$, so $q\in A_{k,\ell}$, and then $\mc{N}(p\restriction_{m+1})\subseteq A_{k,\ell}$, so it's open.

    To show it's dense, given $p\in\sym$ and $m\in\omega$, we need to prove $\mc{N}(p\restriction_{m+1})\cap A_{k,\ell}\neq\emptyset$. Take $q\in\mc{N}(p\restriction_{m+1})$ and extend it randomly until $\ell$. Now, extend $q$ for $r\geq\ell$ such that $q(r)\notin\omega\setminus A$ until $d_n(q(A))<\frac1k$, after this choose $q(r)\in A$ such that $d_n(q(A))>1-\frac1k$. Then $q\in A_{k,\ell}$.
\end{proof}

\begin{definition}{$h$-slalom}{h_slalom}
    Let $h\in\omega^\omega$ with $\lvert h(n)\rvert\geq1$, $\forall n\in\omega$. A function $\phi:\omega\to[\omega]^{<\omega}$ is an \lasereye{$h$-slalom} if $\lvert\phi(n)\rvert\leq h(n)$, $\forall n\in\omega$. We denote by $\Phi_h$ the collection of all $h$-slaloms.
\end{definition}

Now, to prove the other tukey inequality, let's take $\mc{A}_1=(\omega^\omega,\Phi_h,\in^\infty)$, where $f\in^\infty\phi$ if $f(n)\in\phi(n)$ for infinitely many $n$, and use the following result due to Bartoszyński.

\begin{lemma}{(Bartoszyński)}{non_cov_bartoszynski} For any $h\in\omega^\omega$ satisfying $\lvert h(n)\rvert\geq1$, $\forall n\in\omega$, we have
    \begin{align*}
        \textbf{non}(\mc{M}) & = \mf{d}\rp{\omega^\omega,\Phi_h,\in^\infty};\\
        \textbf{cov}(\mc{M}) & = \mf{b}\rp{\omega^\omega,\Phi_h,\in^\infty}.
    \end{align*}
\end{lemma}

It's possible to show that $\rp{\omega^\omega,\Phi_h,\in^\infty}$ and $\rp{\mc{M},2^\omega,\not\ni}$ are not Tukey equivalent, so the approach to prove $\mf{dd}\geq\textbf{non}(\mc{M})$ changes a lot from the proof of \cref{lem:dd_non} and it's way trickier.

\begin{lemma}{(\cite{brech})}{} For any $h\in\omega^\omega$ growing fast enough (as in \cref{eq:slalom} in the sketch below)
    $$\rp{\omega^\omega,\Phi_h,\in^\infty}\leq_T(D,\sym,R)$$
\end{lemma}

\begin{sketch}
    Before giving the formal proof, let's sketch what is the intuitive idea behind. We have some $h\in\omega^\omega$ and we want $\varphi_-:\omega^\omega\to D$ and $\varphi_+:\sym\to\Phi_h$ such that
    $$\varphi_-(g)R p\Rightarrow g\in^\infty\varphi_+(p),\,\forall g\in\omega^\omega,\,\forall p\in\sym$$
    Since $\varphi_-(g)Rp$ means that $d(\varphi_-(g))\neq d(p(\varphi_-(g)))$, we will go for the contrapositive, i.e., assume there is $n_0\in\omega$ such that $g(n)\notin\varphi_+(p)(n)$, $\forall n\geq n_0$ and force that $d(\varphi_-(g))=d(p(\varphi_-(g)))$.

    Now, since we have the freedom to choose $h$, let's choose it so we could prove $d(\varphi_-(g))=0$. If we could take $\varphi_+(p)$ such that $2^n\subseteq\varphi_+(p)(n)$, then, since $g(n)\notin\varphi_+(p)(n)$, $\forall n\geq n_0$, we have $g(n)\geq 2^n$, which implies that, if $\varphi_-(g)$ grows as fast as $g$, then $d(\varphi_-(g))=0$. Indeed, we could try to take $\varphi_-(g)=\{2^n\}_{n\in\omega}$, but then we lost the control we have, because we can no longer use our hypothesis $g(n)\notin\varphi_+(p)(n)$, so we will use a clever definition that depends on $g$:
    
    Let $\varphi_-(g)=\{a_n\}_{n\in\omega}$ be such that $(a_n)$ is any increasing sequence, for $n<n_0$; and $a_{n+1}=g(a_n)$, for $n\geq n_0$. Then, since $g(n)\geq 2^n$, $\forall n\geq n_0$, it implies that
    $$a_{n+1}=g(a_n)\geq2^{a_n}\geq2^n\text{, for }n\geq n_0$$
    so we have that $d(\varphi_-(g))=0$ as desired.
    
    Now, it remains to choose $h\in\omega^\omega$ such that $2^n\in\varphi_+(p)(n)$ and $d(p(\varphi_-(g)))=0$. We want that $p$ sparse the elements of $\varphi_-(g)$ as much as $g$, in particular, it would be good if $p(a_{n+2})\geq a_{n+1}\geq2^n$, for all $n\geq n_0$, in this case $d(p(\varphi_-(g)))=0$. To do this, we need that 
    $$p(a_{n+1})\neq 0,1,\dots,a_n\text{, i.e., }a_{n+1}\neq p^{-1}(0),p^{-1}(1),\dots,p^{-1}(a_n)$$
    So let's choose
    $$\varphi_+(p)=2^n\cup\{p^{-1}(0),p^{-1}(1),\dots,p^{-1}(n)\}$$
    then, $a_{n+1}=g(a_n)\notin\varphi_+(p)(a_n)$, $\forall n\geq n_0$ implies $p(a_{n+1})\geq a_n=g(a_{n-1})\geq2^{n-1}$. So the natural choice for $h$ is any $h\in\omega^\omega$ such that
    \begin{equation}\label{eq:slalom}
        h(n)\geq2^n+n+1=\lvert\varphi_+(p)\rvert
    \end{equation}
\end{sketch}

\begin{proof}
    Let $h\in\omega^\omega$ be such that $h(n)\geq2^n+n+1$, $\forall n\in\omega$. Choose $\varphi_+(p)(n)=2^n\cup\{p^{-1}(k):k\leq n\}$, then clearly $\varphi_+(p)\in\Phi_h$. Now, assume there is some $n_0\in\omega$ such that $g(n)\notin\varphi_+(p)(n)$, and therefore $g(n)\geq2^n$, $\forall n\geq n_0$. Define $\varphi_-(g)=\{a_n:n\in\omega\}$ where $(a_n)$ is any increasing sequence for $n<n_0$, and $a_{n+1}=g(a_n)$, for $n\geq n_0$. Since $a_{n+1}=g(a_n)\notin\varphi_+(p)(a_n)$, $\forall n\geq n_0$, then $a_{n+1}\neq p^{-1}(0)\dots,p^{-1}(a_n)$, i.e., $p(a_{n+1})\neq0,1,\dots,a_n$, so $p(a_{n+1})>a_n\geq2^{n-1}$, which means $d(\varphi_-(g))=0=d(p(\varphi_-(g)))$.
\end{proof}

With such characterization in hands, and given that we have defined $\mf{dd}$ as an analogous cardinal to $\mf{rr}$, it's a natural question to ask if $\mf{rr}=\textbf{non}(\mc{M})$ also. In fact, it's far from clear if this is true. The first reason is that, in some sense, $\mf{dd}$ is not actually the analogous cardinal to $\mf{rr}$, since c.c. series corresponds to sets with density strictly between $0$ and $1$, rather than p.c.c. series diverging to $\pm\infty$. The following table summarizes this correspondence:

\begin{center}
\begin{tabular}{ |c|c|c| } 
\hline
p.c.c. serie $\sum_n a_n$ & infinite-coinfinite set $A$\\
\hline
c.c. series & $d(A)\in(0,1)$ \\ 
$\sum_na_n=\pm\infty$ & $d(A)\in\{0,1\}$ \\ 
$\sum_n a_n$ diverges by oscillation & $d(A)$ is undefined\\
\hline
\end{tabular}
\end{center}

where $\sum_n a_n$ is p.c.c. (potentially conditionally convergent) if there is some $p\in\sym$ such that $\sum_na_{p(n)}$ is c.c. In other words, it's c.c. without the condition that the original series converges to a finite number. With this in mind, we can define the actual cardinal analogous of $\mf{dd}$:

\begin{definition}{$\mf{rr}'$}{rr_prime}
    \lasereye{$\mf{rr}'$} is the smallest cardinality of any family $C\subseteq\sym$ such that, for every p.c.c. series $\sum_n a_n$ of reals, there is some $p\in C$ for which $\sum_n a_{p(n)}$ no longer converges to the same limit.
\end{definition}

using the analogous proofs for $\mf{rr}$, it's not hard to see that $\mf{rr}\leq\mf{rr}'\leq\textbf{non}(\mc{M})$. In fact, we have:

\begin{theorem}{(\cite{brech})}{rr_prime_non}
    $$\mf{rr}'=\textbf{non}(\mc{M})$$
\end{theorem}

whose proof is as the proof for $\mf{dd}$, since it's sufficient to show that $(\omega^\omega,\Phi_h,\in^\infty)\leq_T(\text{p.c.c.},\sym,R)$

\begin{proof}
    Let $\varphi_+$ be the same, i.e., $\varphi_+(p)(n)=2^n\cup\{p^{-1}(k):k\leq n\}$, then, if $g\in\omega^\omega$ is such that $g(n)\notin\varphi_+(p)(n)$ for all $n\geq n_0$, for some $n_0\in\omega$, then, let $A=\{a_n:n\in\omega\}$ (which was $\varphi_-(g)$ in the previous proof) and $B=\{b_0<b_1<\dots\}=\omega\backslash A$. We define $\varphi_-(g)=c_k$, where
    $$c_k=\begin{cases}
        \frac1n\text{, if }k=b_n\\
        -\frac1n\text{, if }k=a_n
    \end{cases}$$
    clearly $\sum_k c_k$ is p.c.c. and, since $a_n\geq2^{n-1}$, we have $\sum_kc_k=\infty$. Also, since by our choice of $h$, we have $p(a_{n+2})\geq2^n$, then $\sum_kc_{p(k)}=\infty$ \todo{why?}for all $p$.
\end{proof}

\section{Variations on \texorpdfstring{$\mf{dd}$}{dd}}

Motivated by the fact that $\mf{rr}'=\textbf{non}(\mc{M})=\mf{dd}$, we will generalize the definition of $\mf{dd}$ as following:

\begin{definition}{$(X,Y)$-density number}{xy_density_number}
    Let $X,Y\subseteq\textbf{all}=[0,1]\cup\{\osc\}$ where $X\neq\emptyset$ and for all $x\in X$ there is some $y\in Y$ such that $y\neq x$. The \lasereye{$(X,Y)$-density number $\mf{dd}_{X,Y}$} is the smallest cardinality of a family $C\subseteq\sym$ such that for every infinite-coinfinite set $A\subseteq\omega$ with $d(A)\in X$, there is some $p\in C$ such that $d(p[A])\in Y$ and $d(p[A])\neq d(A)$.
\end{definition}

The first two conditions are to guarantee that $\mf{dd}_{X,Y}$ is well-defined, since we need $d(A)\in X$ (then $X\neq\emptyset$) and, if $d(A)=x\in X$, we need some $y\in Y$, such that $d(p[A])=y\neq x$.

With this notation, it's easy to see that $\mf{dd}_{[0,1],\textbf{all}}=\mf{dd}$ and, if $X'\subseteq X$ and $Y'\supseteq Y$, then
$$\mf{dd}_{X',Y'}\leq\mf{dd}_{X,Y}$$
Let's now provide the Galois-Tukey framework for $\mf{dd}_{X,Y}$:
$$\mf{d}(D_X,\sym,R_Y)=\mf{dd}_{X,Y}$$
where $D_X$ is the family of infinite-coinfinite sets $A$ with $d(A)\in X$ and $AR_Yp$ if $d(p[A])\in Y$ and $d(p[A])\neq d(A)$, for $A\in D_X$ and $p\in\sym$.

Let's then rephrase the main results of the last section in this framework:

\begin{theorem}{(\cite{brech})}{nonM_general}
    The following hold:
    \begin{enumerate}
        \item If $\osc\notin X$ and $\osc\in Y$, then
        $$(D_X,\sym,R_Y)\leq_T(\mc{M},2^\omega,\not\ni);$$
        \item If $0\in X$ or $1\in X$, then
        $$(\omega^\omega,\Phi_h,\in^\infty)\leq_T(D_X,\sym,R_Y)$$
        for $h\in\omega^\omega$ with $h(n)\geq 2^n+n+1$, for all $n\geq0$.
    \end{enumerate}
\end{theorem}

\begin{proof}
    1. Is easy to see that follows from \cref{lem:dd_non} since, if $\osc\notin X$ and $\osc\in Y$, we can repeat the construction in the proof and find some $p\in\sym$ such that $d(p[A])=\osc\in Y\setminus X$.

    2\todo{idk}.
\end{proof}

This result already characterizes a lot of $\mf{dd}_{X,Y}$ as $\textbf{non}(\mc{M})$. We will see later that, in fact, this result can't be improved in the following sense: the assumptions in 1 and 2 are necessary.

With such characterization, it's natural to look for other bounds of $\mf{dd}_{X,Y}$ for when $0,1\notin X$, $\mathbf{osc}\in X$ or $\mathbf{osc}\notin Y$. We will mostly obtain lower bounds, but there is one upper bound too.

\begin{theorem}{(\cite{brech})}{cov_ddxy}
    If $X\cap[0,1]\neq\emptyset$, then
    $$(2^\omega,\mc{N},\in)\leq_T(D_X,\sym,R_Y)$$
\end{theorem}

Note that the case when $0\in X$ or $1\in X$ can be reduced to the previous theorem if we could prove thar $(2^\omega,\mc{N},\in)\cong(\omega^\omega,\mc{N},\in)\leq_T(\omega^\omega,\Phi_h,\in^\infty)$. In fact, the proof goes by these lines. We will not prove it, since it needs some measure theory concepts.

It's worth noting that it's an open question if the previous theorem holds if $X=\{\osc\}$.

\begin{corollary}{(\cite{brech})}{cov_ddxy}
    If $X\cap[0,1]\neq\emptyset$, then $\textbf{cov}(\mc{N})\leq\mf{dd}_{X,Y}$
\end{corollary}

we can also compare $\mf{dd}_{X,Y}$ with $\mf{b}$ and $\mf{d}$ when $\osc\in X$ or $\osc\notin Y$ as follows

\begin{theorem}{(\cite{brech})}{b_ddxy}
    If $\osc\in X$ or $\osc\notin Y$, then
    $$(\omega^\omega,\omega^\omega,\not\geq^*)\leq_T(D_X,\sym,R_Y)$$
\end{theorem}

This proof is gigantic, so let's break it in some proof-sketches, one for each case.

\begin{sketch}
    Assume first $\text{osc}\in X$, we want $\varphi_-:\omega^\omega\to D_X$ and $\varphi_+:\sym\to\omega^\omega$ such that
    $$\forall f\in\omega^\omega,\forall p\in\sym,\,f\geq^*\varphi_+(p)\implies d(p[\varphi_-(f)])=d(\varphi_-(f))\in Y$$
    Since we only know $\text{osc}\in X$, we want $\varphi_-(f)$ s.t. $d(\varphi_-(f))=\text{osc}$. In fact, following the intuition of \ref{lem:almost_disjoint_osc}, we will construct such disjoint intervals using the growth of $f$, so we first need to choose $\varphi_+(p)$ in a way that $f\geq^*\varphi_+(p)$ guarantees us that we can let $a_n=f(a_{n-1})$ increases fast.

    Let, for now, $\varphi_+(p)(n)\geq 2^n$, then $f(n)\geq 2^n$ eventually, i.e., $\forall n\geq n_0$, let's then define $a_0^f=0$ and
    $$a_n^f=\begin{cases}
        2^{a^f_{n-1}},\text{ if }n<n_0\\
        f(a_{n-1}^f),\text{ if }n\geq n_0
    \end{cases}$$
    we could then just let $\varphi_-(f)=\bigcup_{n\geq0}[a_{2n}^f,a_{2n+1}^f)$, and try to adjust $\varphi_+(p)$ in such a way that we can control where $p(k)$ goes knowing that $k<a_{2n}^f$ or $k>a_{2n+1}^f$ but, to do that in a satisfactory way, it's better to consider a middle point between $a_{2n}^f$ and $a_{2n+1}^f$ such that everyone from the left outside of this interval remains to the left, and similarly for the right.

    Let then $\varphi_-(f)=\bigcup_{n\geq0}[a_{4n}^f,a_{4n+2}^f)$. We want that if $n_0<k<a_{4n}^f$, then $p(k)<a_{4n+1}^f=f(a_{4n+1}^f)$. Then, it would be good if we had $p(k)<\varphi_+(p)(k)$, since that could imply
    $$p(k)<\varphi_+(p)(k)<\varphi_+(p)(a_{4n}^f)\leq f(a_{4n}^f)=a_{4n+1}^f$$
    let then $\varphi_+(p)(n)\geq\max\{p(k):k\leq n\}$. And, since we also want $p^{-1}(k)<a_{4n+2}^f$, whenever $k<a_{4n+1}^f$, we can then finally define
    $$\varphi_+(p)(n)=\max\{p(k),p^{-1}(k):k\leq n\}+2^n$$
\end{sketch}

\begin{proof}
    Assume $\text{osc}\in X$ and let $\varphi_+(p)(n)=\max\{p(k),p^{-1}(k):k\leq n\}+2^n$. Since there is some $n_0\in\omega$ such that $f(n)\geq\varphi_+(p)(n)$, $\forall n\geq n_0$, let $a_0^f=0$ and
    $$a_n^f=\begin{cases}
        2^{a^f_{n-1}},\text{ if }n<n_0\\
        f(a_{n-1}^f),\text{ if }n\geq n_0
    \end{cases}$$
    then, define $\varphi_-(f)=\bigcup_{n\geq0}[a_{4n}^f,a_{4n+2}^f)$. So, as we saw, if $k<a_{4n}^f$, then $p(k)<a_{4n+1}^f$ and, if $k\geq a_{4n+2}^f$, then $p(k)\geq a_{4n+1}^f$; so $p[a_{4n+2}^f]=a_{4n+2}^f$. Now, we want to show that
    $$\lvert p[\varphi_-(f)]\cap a_{4n+1}^f\rvert=\lvert\varphi_-(f)\cap a_{4n+1}^f\rvert$$
    but it's not that hard to see. By the mere definition of $\varphi_-(f)$, the elements before $a_{4n+1}^f$ are the ones until ``half'' of $I_n=[a_{4n}^f,a_{4n+2}^f)$; and the ones that get mapped to elements before $a_{4n+1}^f$ are, as we saw, first all the elements before $I_n$, and then $a_{4n+1}^f-a_{4n}^f$ elements inside $I_n$, resulting in the same amount. Similarly, using $p[a_{4n+2}^f]=a_{4n+2}^f$, we can conclude that
    $$\lvert p[\varphi_-(f)]\cap a_{4n+3}^f\rvert=\lvert\varphi_-(f)\cap a_{4n+3}^f\rvert$$
    but then, we have that
    \begin{align*}
        d_{a_{4n+1}^f}(p[\varphi_-(f)]) & = d_{a_{4n+1}^f}(\varphi_-(f))\geq\frac{a_{4n+1}^f-a_{4n}^f}{a_{4n+1}^f}=1-\frac{a_{4n}^f}{a_{4n+1}^f}\to1\\
        d_{a_{4n+3}^f}(p[\varphi_-(f)]) & = d_{a_{4n+3}^f}(\varphi_-(f))\leq\frac{a_{4n+2}^f}{a_{4n+3}^f}\to0
    \end{align*}
    so $d(p[\varphi_-(f)])=\text{osc}$.
\end{proof}

\begin{sketch}
    Now, let's assume $\text{osc}\notin Y$. In order to illustrate the main idea, let's first consider the case where $\frac12\in X$. Say $p\in\sym$ is big on $E=2\mbb{N}$ if there is some $k=k_p$ such that there are infinitely many $m$ satisfying
    \begin{equation}\label{eq_big_on_e}
        d_m(p[E])>\frac12+\frac1k
    \end{equation}
    Similarly, it's small on $E$ if there is some $k=k_p$ such that for infinitely many $m$'s
    \begin{equation}\label{eq_small_on_e}
        d_m(p[E])<\frac12-\frac1k
    \end{equation}
    The idea is, if our $p$ is big (or small) on $E$, we can, as before, assume $f\geq^*\varphi_+(p)$ and choose it such that we can make a fast growing sequence $a_n=f(a_{n-1})$ and, with this, define
    \begin{equation}\label{eq_phi_-_def}
        \varphi_-(f)=\bigcup_{n\in\omega}([a^f_{2n},a^f_{2n+1})\cap E)\cup\bigcup_{n\in\omega}([a^f_{2n+1},a^f_{2n+2})\cap(\omega\setminus E))
    \end{equation}
    then, intuitively, $d(\varphi_-(f))=\frac12$, but its rearrangement will oscillate, since we are inside $E$ in a big chunk, but then outside of it, so $p$ in $\varphi_-(f)$ will be sometimes small, sometimes big.

    Since $\text{osc}\notin Y$, this would be enough. Lastly, for the case where $p$ is neither big nor small on $E$, we just need to show that $p[\varphi_-(f)]$ has a subsequence of the density converging to $1/2$, then it's either converges to $1/2$ or oscillates. So, in both cases, we won't have $\varphi_-(f)R_Yp$.
\end{sketch}

\begin{proof}
    If $p$ is big (small) on $E$, let
    $$m_n^p=\min\{m\geq 2^n:m\text{ satisfies eq. }(\ref{eq_big_on_e})(\text{eq. }(\ref{eq_small_on_e}))\}$$
    otherwise, if $p$ is neither big nor small, let
    $$m_n^p=\min\left\{m\geq2^n:\frac12-\frac1{2^n}\leq d_m(p[E])\leq\frac12+\frac1{2^n}\right\}$$
    finally, let's define
    $$\varphi_+(p)(n)=\min\{k>m_n^p:p^{-1}[m_n^p]\subseteq k\}$$
    Now, if $f\geq^*\varphi_+(p)$, then there is some $n_0\in\omega$ such that $f(n)\geq\varphi_+(p)(n)\geq m_n^p\geq 2^n$, $\forall n\geq n_0$. Now, define $(a_n)$ as in the previous proof, and let $\varphi_-(f)$ as in (\ref{eq_phi_-_def}).

    First of all, let's prove properly that $d(\varphi_-(p))=\frac12$. Given some $m\geq n_0$, we know $m\in[a_k^f,a_{k+1}^f)$ for some $k\in\omega$. Since $\lvert[a_i^f,a_{i+1}^f)\cap E\rvert=\frac{a_{i+1}^f-a_i^f}{2}+\varepsilon_i$, where the error $\varepsilon_i\in\{0,\pm1\}$ (similarly for $O$), then, we have that
    \begin{align*}
        d_m(\varphi_-(p)) & = \frac1m\rp{\sum_{n=0}^{k-1}\rp{\frac{a_{n+1}^f-a_n^f}{2}+\varepsilon_n}+\rp{\frac{m-a_k^f}{2}+\varepsilon_k}}\\
        & = \frac{a^f_k-\cancelto{0}{a_0^f}}{2m}+\frac{m-a_k^f}{2m}+\frac{\varepsilon}{m}\\
        & = \frac12+\frac{\varepsilon}m
    \end{align*}
    Since $\varepsilon\leq k$, and $m\geq a_k^f\geq 2^k$, then $\varepsilon/m\to0$.

    Now let's divide in two cases. First, assume $p$ is big on $E$ (small is similar) and fix $2n\geq n_0$, then
    $$2^{a_{2n}^f}\leq m_{a_{2n}^f}^p<\varphi_+(p)(a_{2n}^f)\leq f(a_{2n}^f)=a_{2n+1}^f$$
    We will now prove a seemly random inequality that will actually be important later:
    
    Since $a_{2n+1}^f\geq\varphi_+(p)(a_{2n}^f)$, by definition we have $p^{-1}\rp{m_{a_{2n}^f}^p}\subseteq a_{2n+1}^f$, then $m_{a_{2n}^f}^p\subseteq p[a_{2n+1}^f]$. In particular, every number greater than $a_{2n+1}^f$ has image under $p$ greater than $m_{a_{2n}^f}^p$. This implies that every number in $p[E]\cap m_{a_{2n}^f}^p$ is automatically in $p[E\cap a_{2n+1}^f]\cap m_{a_{2n}^f}^p$, so
    \begin{equation}\label{eq_d(p[E])}
        d_{m_{a_{2n}^f}^p}(p[E\cap a_{2n+1}^f])=d_{m_{a_{2n}^f}^p}(p[E])>\frac12+\frac1{k_p}
    \end{equation}
    Now, with this in hands, we can calculate $d_{m_{a_{2n}^f}^p}(p[\varphi_-(f)])$ and $d_{m_{a_{2n+1}^f}^p}(p[\varphi_-(f)])$ and show it oscillates. Fix $a_{2n+1}^f$, then $\varphi_-(f)$ agrees with $E$ from $a_{2n}^f$ to $a_{2n+1}^f$, then $p[\varphi_-(f)]$ and $p[E]$ also, so
    $$\left\lvert p[\varphi_-(f)]\cap m_{a_{2n}^f}^p\right\rvert\geq\left\lvert p[\varphi_-(f)\cap a_{2n+1}^f]\cap m_{a_{2n}^f}^p\right\rvert\geq\left\lvert p[E\cap a_{2n+1}^f]\cap m_{a_{2n}^f}^p\right\rvert-a_{2n}^f$$
    i.e., they disagree by at most $a_{2n}^f$ elements. So, by equation (\ref{eq_d(p[E])}):
    $$d_{m_{a_{2n}^f}^p}(p[\varphi_-(f)])\geq d_{m_{a_{2n}^f}^p}(p[E])-\frac{a_{2n}^f}{m_{a_{2n}^f}^p}>\frac12+\frac1{k_p}-\frac{a_{2n}^f}{2^{a_{2n}^f}}>\frac12+\frac1{k_p}$$
    Similarly, by equation (\ref{eq_d(p[E])}) we have for $p[O]$ a density $<\frac12-\frac1{k_p}$, and the argument follows analogously.

    So we have $d_{m_{a_{2n}^f}^p}(p[\varphi_-(f)])>\frac12+\frac1{k_p}$ and $d_{m_{a_{2n+1}^f}^p}(p[\varphi_-(f)])<\frac12-\frac1{k_p}$, thus $d(p[\varphi_-(f)])=\text{osc}$.

    Let's now assume $p$ is neither big nor small on $E$. The equality in equation (\ref{eq_d(p[E])}) didn't assume $p$ was big or small, so we know that
    $$d_{m_{a_{n}^f}^p}(p[E\cap a_{n+1}^f])\to\frac12$$
    in particular, the same is true for $\omega\setminus E$ instead. Then, within such subsequence, we can conclude that $d_{m_{a_n^f}^p}(p[\varphi_-(f)])\to\frac12$, which means its density is either $1/2$ or oscillates, as desired.
    
    \noindent\rule{16.6cm}{0.4pt}
    
        Let's now analyze the general case. If $0\in X$ or $1\in X$, the result follows by \cref{th:nonM_general}. Let's then assume $X\subseteq(0,1)$. Given $r\in(0,1)$, let:
    \begin{itemize}
        \item $(\ell_b:b\in\omega)$ be a sequence such that $\left\lvert r-\frac{\ell_b}{2^b}\right\rvert<\frac1{2^b}$;
        \item $(j_b:b\in\omega)$ be such that $j_0=0$ and $j_{b+1}=j_b+2^{2b}$ and $J_b=[j_b,j_{b+1})$;
        \item $J_b=I_b^0\;\vert\; I_b^1\;\vert\;\dots\;\vert\;I_b^{2^b-1}$ be a uniform partition of $J_b$ into $2^b$ intervals of length $2^b$;
        \item for any $c\in\omega$, $[2^c]^{\ell_c}=\{e_c(v):v<\omega_c\}$ be an enumeration of $[2^c]^{\ell_c}$, where $\omega_c=\lvert[2^c]^{\ell_c}\rvert=\binom{2^c}{\ell_c}$ and $e_c(0)=\ell_c$;
        \item for any $c\in\omega$ and $b\geq c$, $I_a^b=L_{b,0}^{a,c}\;\vert\;L_{b,1}^{a,c}\;\vert\;\dots\;\vert\;L_{b,2^c-1}^{a,c}$ be a uniform partition of $I_a^b$ into $2^c$ intervals of length $2^{b-c}$;
        \item for $v<\omega_c$, $E_v^c\subseteq\omega$ is such that $E_v^c\cap j_c=\emptyset$ and, for $b\geq c$ and $a\in 2^b$,
        $$E_v^c\cap I_b^a=\bigcup_{d\in e_c(v)}L_{b,d}^{a,c}$$
    \end{itemize}
    \noindent\rule{16.6cm}{0.4pt}

    Let's first show that $E_v^c$ in fact has density $\ell_c/2^c$: First of all, note that from each $I^a_b$ we will choose $\ell_c$ sets $L_{b,d}^{a,c}$, whose size is $2^{b-c}$, then
    $$\lvert E_v^c\cap I^a_b\rvert=\ell_c\cdot 2^{b-c}=\frac{\ell_c}{2^c}\cdot 2^b$$
    which implies that
    \begin{equation}\label{Evc_cap_Jb_size}
        \lvert E_v^c\cap J_b\rvert=\frac{\ell_c}{2^c}\cdot\lvert J_b\rvert=\frac{\ell_c}{2^c}\cdot 2^{2b}
    \end{equation}
    Now, let $n\in\omega$ be large enough, then $n\in J_b$ for some $b$, write $n=j_b+k\cdot2^b+r_0$, where $k$ is the number of $I_a^b$ blocks before $n$, and $r_0$ the remainder for the incomplete block.

    Now, from $0$ to $j_c$, $E_v^c$ has no elements. From $j_c$ to $j_b$, by equation (\ref{Evc_cap_Jb_size}), we have $\frac{\ell_c}{2^c}(j_b-j_c)$ elements. From $j_b$ to $n$ we have $\frac{\ell_c}{2^c}\cdot(k\cdot 2^b)$ and a remainder of $\widetilde{r}\leq r_0$, therefore
    $$d_n(E_v^c)=\frac{\ell_c}{2^c}\rp{\frac{j_b-j_c+k\cdot 2^b}n}+\frac{\widetilde{r}}{n}=\frac{\ell_c}{2^c}\rp{\frac{n-j_c-r_0}n}+\frac{\widetilde{r}}{n}=\frac{\ell_c}{2^c}-\frac{\ell_c}{2^c}\rp{\frac{j_c+r_0}n}+\frac{\widetilde{r}}{n}$$
    now, $\frac{\widetilde{r}}n\leq\frac{2^b}{j_b}\to0$, and, similarly, $\frac{j_c+r_0}n\to0$, so $d_n(E_v^c)\to\frac{\ell_c}{2^c}$.

    \noindent\rule{16.6cm}{0.4pt}

    Now, as in the previous case, we call $p\in\sym$ big if there is some $k=k_p$ such that for infinitely many $c$ and $m$, we have $d_m(p[E_0^c])>\frac{\ell_c}{2^c}+\frac1k$. Analogously for small and neither. We let $m_n^p$, $\varphi_+(p)$ and $a^f_n$, where $f\geq^*\varphi_+(p)$, as before.

    We then define $(K_c:c\in\omega)$ an interval partition of $\omega$ into intervals of length $\omega_c$. Now, if $n$ is the $v$-th element of $K_c$, for $a_n^f\leq b<a_{n+1}^f$ and $a\in2^b$, define
    $$\varphi_-(f)\cap I_b^a=\bigcup_{d\in e_c(v)}L_{b,d}^{a,c}$$
    Now, in order to show that $d(\varphi_-(f))=r$, we follow a similar approach as $E_v^c$. Let $m\in\omega$, then $m\in J_b$ for some $b$. Also, for each $J_b$, with $b\in[a_n^f,a_{n+1}^f)$, we have a unique $K_c$ and $v$ associated, in particular, we can associate each $J_b$ with some $E_{v(b)}^{c(b)}$ such that $\varphi_-(f)\cap[a_n^f,a_{n+1}^f)=E_{v(b)}^{c(b)}\cap[a_n^f,a_{n+1}^f)$.

    Let then $m=j_b+k\cdot 2^b+r_0$ as before. Let's now use some asymptotic notation (because it's easier, since we already know the argument). Since $\frac{\ell_c}{2^c}\to r$, then from some $c(b_0)$ large enough, we have:
    \begin{align*}
        \lvert\varphi_-(f)\cap m\rvert & = \sum_{i=0}^{b-1}\lvert\varphi_-(f)\cap J_i\rvert+k\lvert\varphi_-(f)\cap I_b^a\rvert+\widetilde{r}\\
        & =\sum_{i=0}^{b-1}\frac{\ell_{c(i)}}{2^{c(i)}}2^{2i}+k\frac{\ell_{c(b)}}{2^{c(b)}}2^b+\widetilde{r}\\
        & \sim C+\sum_{i=b_0}^{b-1}r2^{2i}+kr2^b+\widetilde{r}\\
        & = C+r(m-j_{b_0}-r_0)+\widetilde{r}
    \end{align*}
    then, dividing by $m$, we finally have that
    $$d_m(\varphi_-(f))\sim\cancelto0{\frac{C}{m}}+r\rp{1-\cancelto0{\frac{j_{b_0}-r_0}{m}}}+\cancelto0{\frac{\widetilde{r}}{m}}\to r$$

    \noindent\rule{16.6cm}{0.4pt}

    Let's now first note two important things:
    \begin{enumerate}
        \item We know by definition that $m=m_{a_n^f}^p<\varphi_+(p)(a_n^f)\leq f(a_n^f)=a_{n+1}^f$, then $p^{-1}[m]\subseteq a_{n+1}^f$, so
        \begin{align*}
            \lvert p[\varphi_-(f)]\cap m\rvert & = \lvert p[\varphi_-(f)\cap a_{n+1}^f]\cap m\rvert\\
            \lvert p[E_v^c]\cap m\rvert & = \lvert p[E_v^c\cap a_{n+1}^f]\cap m\rvert
        \end{align*}
        this implies that, since $\varphi_-(f)\cap [a_n^f,a_{n+1}^f)=E_v^c\cap[a_n^f,a_{n+1}^f)$ they differ only at $[0,a_n^f)$, so
        $$\frac{\lvert p[\varphi_-(f)]\cap m\rvert-\lvert p[E_v^c]\cap m\rvert}{m}\leq\frac{a_n^f}{m}=\frac{a_n^f}{m_{a_n^f}^p}\to0$$
        then $d_{m_{a_n^f}^p}(p[\varphi_-(f)])\sim d_{m_{a_n^f}^p}(p[E_v^c])$, where $E_v^c$ is chosen in function of $n$.
        \item Since $d(E_v^c)=\frac{\ell_c}{2^c}$, if $d_m(E_0^c)>\frac{\ell_c}{2^c}+\frac1k$, then we can find another subset $e_c(v)$ instead of $e_c(0)$ such that its elements compensate and $d_m(E_v^c)<\frac{\ell_c}{2^c}$.
    \end{enumerate}
    We can now break in cases:

    \textbf{Case 1. $p$ is big:} The case where $p$ is small is analogous. In this case, let $n(c)\geq n_0$ be the $0$-th item of $K_c$, then its corresponding set is $E_0^c$, so, by (1), we have that
    $$d_{m_{a_{n(c)}^f}^p}(p[\varphi_-(f)])\sim d_{m_{a_{n(c)}^f}^p}(p[E_0^c])>\frac{\ell_c}{2^c}+\frac1{k_p}$$
    but, by (2), if $n'(c)$ is the $v(c)$-th element of $K_c$, we have that
    $$d_{m_{a_{n'(c)}^f}^p}(p[\varphi_-(f)])\sim d_{m_{a_{n'(c)}^f}^p}(p[E_{v(c)}^c])<\frac{\ell_c}{2^c}$$
    so $d(\varphi_-(f))=\text{osc}$, as desired.

    \textbf{Case 2. $p$ is neither big nor small:} then there isn't such $k_p$ as in the definition, which means we can find a sequence $k_c\to\infty$ such that
    $$\frac{\ell_c}{2^c}-\frac1k_c\leq d_{m_{a_{n(c)}^f}^p}(p[\varphi_-(f)])\sim d_{m_{a_{n(c)}^f}^p}(p[E_0^c])\leq\frac{\ell_c}{2^c}+\frac1{k_c}$$
    which implies that, when $c\to\infty$, that $d_{m_{a_{n(c)}^f}^p}(p[\varphi_-(f)])\to r$, so it either oscillates, or converges to $r$, as desired.
\end{proof}

\begin{corollary}{(\cite{brech})}{b_ddxy}
    If $\osc\in X$ or $\osc\notin Y$, then $\mf{b}\leq\mf{dd}_{X,Y}$
\end{corollary}


Now, Consider the relational system $(\mf{c},[\mf{c}]^{\aleph_0},\in)$, let's first prove some properties about it.

\begin{lemma}{(\cite{brech})}{simple_relational_system}
    $\mf{d}(\mf{c},[\mf{c}]^{\aleph_0},\in)=\mf{c}$ and $\mf{b}(\mf{c},[\mf{c}]^{\aleph_0},\in)=\aleph_1$.
\end{lemma}

\begin{proof}
    For the first equality, note that
    $$\mf{d}(\mf{c},[\mf{c}]^{\aleph_0},\in)\leq\lvert[\mf{c}]^{\aleph_0}\rvert=\mf{c}^{\aleph_0}=\rp{2^{\aleph_0}}^{\aleph_0}=2^{\aleph_0\cdot\aleph_0}=2^\aleph_0=\mf{c}$$
    Now, given $A=\{A_\beta\}_{\beta<\kappa}\subseteq[\mf{c}]^{\aleph_0}$ with $\kappa<\mf{c}$, if $U=\bigcup A$, then
    $$\lvert U\rvert\leq\sum_{\beta<\kappa}\lvert A_i\rvert\leq\sum_{\beta<\kappa}\aleph_0=\kappa\cdot\aleph_0<\mf{c}$$
    and, therefore, there is a $\alpha<\mf{c}$ such that $\alpha\notin U$.

    For the second equality, to show it's uncountable, let $X\subseteq\mf{c}$ and just take $Y=X\in[\mf{c}]^{\aleph_0}$, then, for every $\alpha\in X$, we have $\alpha\in Y$ too. Now, given any $X\subseteq\mf{c}$ with $\lvert X\rvert=\aleph_1$, if $Y\in[\mf{c}]^{\aleph_0}$, since it's countable, there is some $\alpha\in X$ such that $\alpha\notin Y$.
\end{proof}

In order to prove the next theorem, we also need the following lemmas about almost disjoint families

\begin{lemma}{(\cite{brech})}{almost_disjoint}
    Let $A\subseteq\omega$ be any infinite set, then there is an almost disjoint family $\{A_\alpha:\alpha<\mf{c}\}$ of subsets of $A$, where almost disjoint means that for any $\alpha\neq\beta<\mf{c}$, we have that $\lvert A_\alpha\cap A_\beta\rvert<\aleph_0$.
\end{lemma}

\begin{proof}
    Since $A\subseteq\omega$ is infinite, there is a bijection $f:\mbb{Q}\to A$, then it's sufficient to find an almost disjoint family $\{A_\alpha:\alpha<\mf{c}\}$ of subsets of $\mbb{Q}$. Then, for each irrational number $\alpha$, let $\{q_n^\alpha\}_{n\in\omega}\subseteq\mbb{Q}$ be a sequence converging to $\alpha$, and let $A_\alpha=\{q^\alpha_n:n\in\omega\}$. Then, since $\alpha$ is irrational, each $A_\alpha$ is infinite and, if $\alpha\neq\beta$, $q_n^\alpha=q_n^\beta$ only for finitely many $n$'s, so $\{f(A_\alpha):\alpha<\mf{c}\}$ is our family.
\end{proof}

This one is a little more technical and less general than the previous one, but we will use it to prove the next theorem

\begin{lemma}{(\cite{brech})}{almost_disjoint_osc}
    There is an almost disjoint family $\{A_\alpha:\alpha<\mf{c}\}$ of subsets of $\omega$ with density $=\text{osc}$.
\end{lemma}

\begin{proof}
    Let $\{B_\alpha:\alpha<\mf{c}\}$ be an almost disjoint family, guaranteed to exists by the previous lemma. Let's then define $A_\alpha$, for each $\alpha<\mf{c}$ as follows:
    $$A_\alpha:=\bigcup_{n\in B_\alpha}P_n\text{, where }P_n=[(2n)!,(2n+1)!)\cap\omega$$
    Clearly $\{P_n\}_{n\in\omega}$ are disjoint and, if $\alpha\neq\beta$, then, by definition, $B_\alpha\cap B_\beta=\{n_1,n_2,\dots,n_m\}$ is finite, and $A_\alpha\cap A_\beta=\cup_{k=1}^mP_k$, which is also finite. Now, let $B_\alpha=\{b_0,b_1,\dots\}$, then, for $N=(2b_n+1)!$, we have $P_{b_n}\subseteq A_\alpha\cap N$, which implies that
    $$d_N(A_\alpha)=\frac{\lvert A_\alpha\cap N\rvert}N\geq\frac{(2b_n+1)!-(2b_n)!}{(2b_n+1)!}=1-\frac{(2b_n)!}{(2b_n+1)!}\to 1$$
    However, if we let $N=(2b_n)!$, then we have that
    $$d_N(A_\alpha)=\frac{\lvert A_\alpha\cap N\rvert}{N}\leq\frac{(2b_n-1)!}{(2b_n)!}\to0$$
    So $d(A_\alpha)=\text{osc}$, as desired.
\end{proof}

\begin{theorem}{(\cite{brech})}{ddxy_c}
    If ($0\in X$ or $1\in X$) and $\text{osc}\notin Y$, or if $\text{osc}\in X$ and ($0\notin Y$ or $1\notin Y$), then
    $$(\mf{c},[\mf{c}]^{\aleph_0},\in)\leq_T(D_X,\sym,R_Y)$$
\end{theorem}

\begin{sketch}
    We need a $\varphi_-:\mf{c}\to D_X$, $\alpha\mapsto A_\alpha$ and a $\varphi_+:\sym\to[\mf{c}]^{\aleph_0}$ such that
    $$\forall\alpha<\mf{c},\forall p\in\sym,\;d(A_\alpha)\neq d(p[A_\alpha])\in Y\Rightarrow\alpha\in\varphi_+(p)$$
    It would be great if we could take $\varphi_+(p)=\{\alpha:d(A_\alpha)\neq d(p[A_\alpha])\in Y\}$, but we don't know if such $\varphi_+(p)$ is in $[\mf{c}]^{\aleph_0}$. Since we don't even know what is in $Y$, let's first consider the case where $0\in X$, then it's sufficient to take $\varphi_+(p)=\{\alpha:d(p[A_\alpha])>0\}$, whenever $\text{osc}\notin Y$. Now, how many sets there are in $\varphi_+(p)$? Well, we can write it as
    $$\varphi_+(p)=\bigcup_{k\geq 1}\underbrace{\left\{\alpha:d(p[A_\alpha])\geq\frac1k\right\}}_{B_k}$$
    And, if we could guarantee that $d(p[A_\alpha\cup A_\beta])=d(p[A_\alpha])+d(p[A_\beta])$, then
    $$\sum_{\alpha\in B_k}d(p[A_\alpha])\geq\frac{\lvert B_k\rvert}{k}$$
    but $\lvert B_k\rvert/k\leq1$, otherwise $p[\bigcup_{\alpha\in B_k}A_\alpha]$ would have density $>1$, then each $\lvert B_k\rvert\leq k$ is finite.

    The easiest scenario would then be where $A_\alpha$ are all pairwsise disjoint, but can we write any infinite $A\subseteq\omega$ as a disjoint family of $\mf{c}$ subsets of it? Of course you can't, otherwise you could select a different element from each subset, and then forming a subset of $A$ with $\mf{c}$ many elements.

    This is where our previous lemma comes in, if $A$ and $B$ are almost disjoint, since density is blind for finitely many sets, then $d(A\cup B)=d(A)+d(B)$.
\end{sketch}

\begin{proof}
    For the first case, assume that $0\in X$ and $\text{osc}\notin Y$. Let $A$ be such that $d(A)=0$, and let $\{\varphi_-(\alpha)=A_\alpha:\alpha<\mf{c}\}$ be an almost disjoint family of subsets of $A$. Then $d(A_\alpha)=0$ for all $\alpha<\mf{c}$ and we can define $\varphi_-(\alpha)=A_\alpha$. Now, given $p\in\sym$, let $\varphi_+(p)=\{\alpha<\mf{c}:d(p[A_\alpha])>0\}$, as we saw this is at most countable, so if it's finite we can just fill in with a bunch of other random $\beta<\mf{c}$ to make it countable. The case where $1\in X$ just consider $\varphi_-(\alpha)=\omega\setminus A_\alpha$, then $d(p[\omega\setminus A_\alpha])=1-d(p[A_\alpha])$ and $\varphi_+(p)=\{\alpha<\mf{c}:d(p[\varphi_-(\alpha)])<1\}=\{\alpha<\mf{c}:d(p[A_\alpha])>0\}$.

    For the second case, let $0\notin Y$ (when $1\notin Y$ just take complements as in $1\in X$). Let, by \cref{lem:almost_disjoint_osc}, $\{\varphi_-(\alpha):\alpha<\mf{c}\}$ be an almost disjoint family of subsets of $\omega$ without density. Then
    $$\varphi_+(p)=\{\alpha<\mf{c}:d(p[\varphi_-(\alpha)])>0\}$$
    is at most countable and satisfies what we want.
\end{proof}

\begin{corollary}{(\cite{brech})}{ddxy_c}
    Under the assumptions of the previous theorem
    $$\mf{dd}_{X,Y}=\mf{c}\text{ and }\mf{dd}_{X,Y}^\bot\leq\aleph_1$$
\end{corollary}

In fact, we can't prove $\mf{dd}_{X,Y}^\bot=\aleph_1$. For example, let $X=\{0\}$ and $Y=(\tfrac12,1]$, then given $F=\{A,B\}$ two disjoint infinite sets of density $0$, if $d(p[A])\in Y$, i.e., $d(p[A])>\frac12$, then $d(p[B])\leq\frac12$, so $B$ witnesses $\neg BR_Yp$ and, therefore, $\mf{dd}_{\{0\},(1/2,1]}^\bot=2$. Analogously, let $Y_n=(\frac1n,1]$, then, if $\{A_k:k\leq n\}$ is any disjoint family of infinite subsets of density $0$, then, we can't have $d(p[A_k])>\frac1n$ for all $k\leq n$, otherwise $p[A_1\cup\dots\cup A_n]$ would have density $>1$. Then $\mf{dd}_{\{0\},Y_n}^\bot=n$.

Let $UR$ be the unreaping relation on $[\omega]^\omega$, where, $A\;UR\;B$ iff $B\subseteq^*A$ or $B\cap A$ is finite (i.e., $A$ does not split $B$ in infinitely many sets inside and outside $A$). $\mf{d}([\omega]^\omega,[\omega]^\omega,UR)=\mf{r}$ is called the reaping number, and $\mf{b}([\omega]^\omega,[\omega]^\omega,UR)=\mf{s}$ is called the splitting number.

\begin{theorem}{(\cite{brech})}{ddxy_r_s}
    Assume $0,1\notin X$ and either
    \begin{itemize}
        \item $0,1\in Y$ or
        \item $\text{osc}\notin X$ and $\exists r\in(0,1]$ s.t. $\forall x\in X$ both $r+x(1-r)$ and $x(1-r)$ are in $Y$.
    \end{itemize}
    Then $(D_X,\sym,R_Y)\leq_T([\omega]^\omega,[\omega]^\omega,UR)$.
\end{theorem}

\begin{sketch}
    We need $\varphi_-:D_X\to[\omega]^\omega$ and $\varphi_+:[\omega]^\omega\to\sym$ such that
    $$y\;UR\;\varphi_-(x)\Rightarrow \varphi_+(y)R_Yx$$
    Let's start with the first condition, where $0,1\notin X$ and $0,1\in Y$. The easiest $\varphi_-$ we could take is the identity, then we need for each $y\in[\omega]^\omega$ a $p_y\in\sym$ such that
    $$x\setminus y\text{ or }x\cap y\text{ is finite}\Rightarrow d(x)\neq d(p_y[x])\in Y$$
    In fact, since we only know $0,1\in Y$, we want $p_y$ such that $d(p_y[x])=0$ or $1$, and the fact that $0,1\notin X$ will already guarantee that $d(x)\neq d(p_y[x])$. Now, $x\setminus y$ being finite, with $x,y\in[\omega]^\omega$, basically says that $y$ is almost inside $x$, except for finitely many elements. And $x\cap y$ being finite says $y$ is almost outside $x$, except for finitely many elements. In both cases, since that finite amount does not affect the density in the limit, then we can see as if $y\subseteq x$ or $y\cap x=\emptyset$.

    In particular, this shows we just need $a,1-a\notin X$ and $a,1-a\in Y$ for any $a\in[0,1]$.
\end{sketch}
    
\begin{proof}
    \begin{itemize}
        \item Let $0,1\in Y$, $\varphi_-(x)=x$ and $\varphi_+(y)=p_y\in\sym$ be such that $d(p_y[y])=1$ for all $y\in[\omega]^\omega$. If $y\subseteq ^*x$, then $F=y\setminus x$ is finite, and $y\subseteq x\cup F$, so
        $$d_n(p_y[x]\cup p_y[F])=\frac{\lvert p_y[x]\cap n\rvert+\lvert p_y[F]\cap n\rvert}n\geq\frac{\lvert p_y[y]\cap n\rvert}n=d_n(p_y[y])$$
        and, since $p_y[F]$ is finite, then the LHS converges to $d(p_y[x])$ and the RHS to $d(p_y[y])=1$.
        
        If $F=x\cap y=x\setminus(\omega\setminus y)$ is finite, then, by the previous argument, $d(p_y[x])\to d(p_y[\omega\setminus y])=0$.

        Now, since $0,1\notin X$, then $d(x)\neq0,1$ and, therefore, $d(x)\neq d(p_y[x])$.
        \item Let $0,1,\text{osc}\notin X$, i.e., $X\subseteq(0,1)$, and assume we have such $r\in(0,1]$. Again, let $\varphi_-(x)=x$ and fix $z\subseteq\omega$ with $d(z)=r$. For $y\in[\omega]^\omega$, let $y'\subseteq y$ be s.t. $d(y')=0$ and let $\varphi_+(y)=p_y$ be the unique permutation mapping $y'$ to $z$ and $\omega\setminus y'$ to $\omega\setminus z$, both in an order-preserving way.

        Now, since $x=(x\cap y')\sqcup(x\setminus y')$, then
        \begin{equation}
            d(p_y[x])=d(p_y[x\cap y'])+d(p_y[x\setminus y'])
        \end{equation}
        Let $\omega\setminus y'=\{a_0<a_1<\dots\}$ and $\omega\setminus z=\{b_0<b_1<\dots\}$, then $p_y(a_i)=b_i$ and, since $x\setminus y'\subseteq\omega\setminus y'$, we will have that, if $(\omega\setminus y')\cap n$ has $k$ elements $a_0,a_1,\dots,a_{k-1}$, then
        $$d_{\omega\setminus y'}^n(x\setminus y')=\frac{\lvert\{i<k:a_i\in x\setminus y'\}\rvert}k$$
        and, as $n\to\infty$, $k\to\infty$, so $d_{\omega\setminus y'}^n(x\setminus y')\to d(x\setminus y')$, but $x=(x\setminus y')\sqcup(y'\cap x)$, where $y'\cap x\subseteq y'$, then $d(y'\cap x)=0$ and, therefore, $d_{\omega\setminus y'}(x\setminus y')=d(x)$.

        Now, since $p_y[\omega\setminus y']=\omega\setminus z$, then, if $(\omega\setminus y')\cap n$ has $k$ elements, then also has $(\omega\setminus z)\cap n$ and, by the same argument as before, we then have that
        $$d_{\omega\setminus z}^n(p_y[x\setminus y'])=\frac{\lvert i<k:b_i\in p_y[x\setminus y']\rvert}{k}=\frac{\lvert\{i<k:a_i\in x\setminus y'\}\rvert}{k}\to d(x)$$
        and, finally, we can then calculate one of the terms in (6), i.e.
        $$d_n(p_y[x\setminus y'])=\frac{\lvert p_y[x\setminus y']\cap n\rvert}{n}=\frac{\lvert p_y[x\setminus y']\cap n\rvert}{\lvert(\omega\setminus z)\cap n\rvert}\cdot\frac{\lvert(\omega\setminus z)\cap n\rvert}{n}\to d(x)(1-r)$$
        Let's now, again, split in both cases: If $y\subseteq^* x$, then $y'\subseteq^*x$, $y'=(x\cap y')\sqcup(y'\setminus x)$ and, since $p_y[y']=z$ and $y'\setminus x$ is finite, then also is $p_y[y'\setminus x]$, and $d(p_y[y'\setminus x])=0$, therefore
        $$r=d(z)=d(p_y[x\cap y'])+d(p_y[y'\setminus x])=d(p_y[x\cap y'])$$
        So, from (6), $d(p_y[x])=d(x)(1-r)+r>d(x)$, since $1\notin X$.

        Now, if $x\cap y$ is finite, then also is $x\cap y'$, therefore $d(p_y[x])=d(x)(1-r)<d(x)$, since $0\notin X$.
    \end{itemize}
\end{proof}

\begin{corollary}{(\cite{brech})}{ddxy_r_s}
    Under the assumptions of the previous theorem
    $$\mf{dd}_{X,Y}\leq\mf{r}\text{ and }\mf{s}\leq\mf{dd}_{X,Y}^\bot$$
\end{corollary}

Finally, we can also prove the following.

\begin{theorem}{(\cite{brech})}{dd_osc_all}
    For any $X\subseteq[0,1]$, we have that
    $$\mf{dd}_{X,\{\text{osc}\}}=\mf{dd}_{X,\text{all}}$$
\end{theorem}

\begin{proof}
    Since $\{\text{osc}\}\subseteq\text{all}$, then $\mf{dd}_{X,\text{all}}\leq\mf{dd}_{X,\{\text{osc}\}}$, so we just need to show that $\mf{dd}_{X,\{\text{osc}\}}\leq\mf{dd}_{X,\text{all}}$.

    Let $C\subseteq\sym$ be a witness for $\mf{dd}_{X,\text{all}}$, and let $\widetilde{C}=C\cup\{g_p:p\in C\}$, where $g_p$ is the function that alternates between $\text{id}$ and $p$ in The Mix Lemma. Now, given $A\in D_X$, let $p\in C$ be such that $d(p[A])\neq d(A)$. If $d(p[A])=\text{osc}$, we are done. Otherwise, we have $d(p[A])\in[0,1]$ and, therefore, $d(g_p[A])=\text{osc}$, which concludes the proof.
\end{proof}